% ===================================================================
%  ਸਾਰਣੀ ਨੰ. 14 — ਸੜਕ। — ਦੁਕਾਨਦਾਰ।
%  Traduction 2026 d'apres texto/io, controlee sur texto/fr, texto/en
%  et texto/hi. Consigne : voir texto/pa/10-sarni-01.tex.
% ===================================================================

%%K t14-apar-1 apar
\VUsaut{4.0mm}
\VUtitre{15.0pt}{60}{ਸਾਰਣੀ ਨੰ. 14}
\VUsaut{3.0mm}

%%K t14-tit-1 sub
\VUpk{11.4pt}{40}{ਸੜਕ। --- ਦੁਕਾਨਦਾਰ।}
\VUsaut{3.0mm}

%%K t14-01-1 p
1. --- ਮੇਰਾ ਮਿੱਤਰ ਜਾਨ, ਜੋ ਪਿੰਡ ਵਿਚ ਰਹਿੰਦਾ ਹੈ, ਮੇਰੇ ਪਰਿਵਾਰ ਨਾਲ ਤਿੰਨ ਦਿਨ ਲੰਘਾਉਣ
ਆਇਆ ਹੈ। \VUgras{ਹੁਣ ਤਕ} ਉਸ ਨੂੰ ਕਦੇ ਕੋਈ ਵੱਡਾ ਸ਼ਹਿਰ ਵੇਖਣ ਦਾ ਮੌਕਾ ਨਾ ਮਿਲਿਆ ਸੀ, ਸੋ
ਅਸੀਂ ਆਪਣੇ ਸਾਰੇ ਦਿਨ ਘੁੰਮਦੇ ਲੰਘਾਉਂਦੇ ਹਾਂ ਤੇ ਜੋ ਉਸ ਨੂੰ ਪਤਾ ਨਹੀਂ, ਉਹ ਮੈਂ ਉਸ ਨੂੰ
ਸਮਝਾਉਂਦਾ ਹਾਂ।

%%K t14-02-1 p
2. --- ਪਹਿਲਾਂ ਅਸੀਂ \VUgras{ਰਸਦਗਾਹ} \textsuperscript{(1)} ਵੇਖਣ ਗਏ, ਜਿਸ ਵਿਚ ਉਸ ਨੂੰ
ਬੜੀ ਦਿਲਚਸਪੀ ਹੋਈ। ਉਸ \VUgras{ਸੜਕ} \textsuperscript{(3)} ਤੋਂ ਮੁੜਦੇ ਹੋਏ ਜਿੱਥੇ
\VUgras{ਤੁਰਕੀ ਹਮਾਮ} \textsuperscript{(2)} ਹਨ, ਸਾਨੂੰ ਇਕ \VUgras{ਸਸਕਾਰ}
\textsuperscript{(4)} ਮਿਲਿਆ। \VUgras{ਸਸਕਾਰ ਦੇ ਜਲੂਸ} ਦੇ ਅੱਗੇ ਅੱਗੇ
\VUgras{ਪੁਰੋਹਿਤ} ਤੁਰ ਰਹੇ ਸਨ, ਉਸ \VUgras{ਮੁਰਦਾ ਗੱਡੀ} ਦੇ ਅੱਗੇ ਜਿਸ ਵਿਚ
\VUgras{ਤਾਬੂਤ} ਰੱਖਿਆ ਸੀ। ਬਹੁਤ ਸਾਰੇ ਮਿੱਤਰ \VUgras{ਸੋਗ} ਵਿਚ ਡੁੱਬੇ ਪਰਿਵਾਰ ਦੇ ਪਿੱਛੇ
ਤੁਰ ਰਹੇ ਸਨ।

%%K t14-02-2 p
ਜਲੂਸ ਦੇ ਲੰਘ ਜਾਣ ਉੱਤੇ ਅਸੀਂ ਵੱਡੇ \VUgras{ਸ਼ਰਾਬਖ਼ਾਨੇ} \textsuperscript{(5)} ਦੇ ਸਾਹਮਣੇ
ਸੜਕ ਲੰਘੀ, ਜਿੱਥੇ ਬਹੁਤ ਸਾਰੇ \VUgras{ਗਾਹਕ} ਛੱਤ ਦੇ ਹੇਠਾਂ \VUgras{ਬੀਅਰ} ਪੀ ਰਹੇ ਸਨ,
ਸ਼ੀਸ਼ੇ ਦੀ \VUgras{ਛਤਰੀ} \textsuperscript{(6)} ਦੀ ਓਟ ਵਿਚ।

%%K t14-03-1 p
3. --- ਜਾਨ ਉਸ \VUgras{ਸ਼ਾਹੀ ਨਿਸ਼ਾਨ} ਤੋਂ ਉਲਝਣ ਵਿਚ ਪੈ ਗਿਆ ਜੋ \VUgras{ਸ਼ਰਾਬ ਦੇ
ਵਪਾਰੀ} \textsuperscript{(7)} ਦੀ ਦੁਕਾਨ ਤੇ \VUgras{ਸਾਜ਼ ਵੇਚਣ ਵਾਲੇ}
\textsuperscript{(10)} ਦੀ ਦੁਕਾਨ ਵਿਚਕਾਰ ਲੱਗਾ ਹੈ। ਉਸ ਨੇ ਮੇਰੇ ਕੋਲੋਂ ਪੁੱਛਿਆ ਕਿ ਉਸ ਦਾ ਕੀ
ਅਰਥ ਹੈ; ਮੈਂ ਆਖਿਆ ਕਿ ਉਹ ਅੰਗਰੇਜ਼ੀ \VUgras{ਕੌਂਸਲਖ਼ਾਨੇ} \textsuperscript{(8)} ਦੀ ਪਛਾਣ
ਹੈ, ਤੇ ਉਸ ਨੂੰ ਬਾਲਾਖ਼ਾਨੇ ਉੱਤੇ ਖੜੇ \VUgras{ਕੌਂਸਲ} \textsuperscript{(9)} ਵੱਲ ਇਸ਼ਾਰਾ
ਕੀਤਾ।

%%K t14-tit-2 sub
\VUpk{10.2pt}{40}{ਦੁਕਾਨਾਂ ਵੇਖਣਾ।}
\VUsaut{3.0mm}

%%K t14-04-1 p
4. --- ਸੁਭਾਵਕ ਹੀ ਅਸੀਂ \VUgras{ਸਾਜ਼ਾਂ}, \VUgras{ਹਾਰਮੋਨੀਅਮਾਂ},
\VUgras{ਆਰਗਨਾਂ} ਤੇ \VUgras{ਪਿਆਨੋ} ਦੀ ਸਜਾਵਟ ਦੇ ਸਾਹਮਣੇ ਰੁਕ ਗਏ; ਅਸੀਂ
\VUgras{ਸੁਰ-ਲਿਪੀਆਂ} ਤੇ ਪਿਆਨੋ ਦੇ \VUgras{ਟੁਕੜਿਆਂ} ਦੇ ਰੰਗਦਾਰ ਸਰਵਰਕ ਵੇਖੇ, ਤੇ ਦੁਕਾਨ
ਦੇ ਬੋਰਡ ਦੇ ਤੌਰ ਉੱਤੇ ਲੱਗੀ ਸੋਨੇ ਦੀ ਲੱਕੜ ਦੀ \VUgras{ਬੀਣਾ} ਤੱਕੀ।

%%K t14-05-1 p
5. --- \VUgras{ਝਾੜੂ ਲਾਉਣ ਵਾਲਿਆਂ} \textsuperscript{(11)} ਤੇ
\VUgras{ਝਾੜੂ ਦੀ ਗੱਡੀ} \textsuperscript{(12)} ਤੋਂ ਉੱਠੇ ਧੂੜ ਦੇ ਬੱਦਲਾਂ ਨੇ ਸਾਨੂੰ ਮਜਬੂਰ
ਕੀਤਾ ਕਿ ਅਸੀਂ \VUgras{ਸਾਮਾਨ ਦੇ ਵਪਾਰੀ} \textsuperscript{(13)} ਦੇ ਸੋਹਣੇ
\VUgras{ਸਾਮਾਨ ਦੇ ਸੈੱਟਾਂ} ਤੇ \VUgras{ਲੋਹੇ ਦੇ ਵਪਾਰੀ} \textsuperscript{(14)} ਦੀ
\VUgras{ਚੁੱਲ੍ਹਿਆਂ} ਤੇ \VUgras{ਲੈਂਪਾਂ} ਦੀ ਸਜਾਵਟ ਕੋਲੋਂ ਛੇਤੀ ਨਿਕਲ ਜਾਈਏ।

%%K t14-06-1 p
6. --- ਉਸੇ ਥਾਂ ਸਾਨੂੰ ਪਟੜੀ ਛੱਡਣੀ ਪਈ, ਕਿਉਂਕਿ ਉਹ ਉਨ੍ਹਾਂ ਗੱਠੜੀਆਂ ਨਾਲ ਰੁਕੀ ਸੀ ਜੋ ਇਕ
\VUgras{ਸਾਮਾਨ ਢੋਣ ਦੀ ਗੱਡੀ} \textsuperscript{(16)} ਤੋਂ ਲਾਹੀਆਂ ਜਾ ਰਹੀਆਂ ਸਨ, ਤਾਂ ਜੋ ਉਸ
\VUgras{ਦੁਕਾਨ} \textsuperscript{(15)} ਵਿਚ ਲੈ ਜਾਈਆਂ ਜਾਣ ਜੋ ਹੁਣੇ ਹੁਣੇ ਕਿਰਾਏ ਉੱਤੇ ਚੜ੍ਹੀ
ਸੀ।

%%K t14-06-2 p
ਇਸ ਨਾਲ ਅਸੀਂ \VUgras{ਗੱਡੀ ਬਣਾਉਣ ਵਾਲੇ} \textsuperscript{(17)} ਦੀਆਂ ਸੋਹਣੀਆਂ ਗੱਡੀਆਂ ਨਾ
ਵੇਖ ਸਕੇ, ਪਰ ਅਸੀਂ ਰੁਕ ਕੇ \VUgras{ਬੰਨ੍ਹਣ ਵਾਲੇ} \textsuperscript{(19)} ਨੂੰ ਇਹ ਵੇਖਣੋਂ
ਨਾ ਰੁਕੇ ਕਿ ਉਹ ਆਪਣੇ ਗੁਆਂਢੀ \VUgras{ਸਾਈਕਲ ਤੇ ਮੋਟਰ ਦੇ ਵਪਾਰੀ} \textsuperscript{(18)}
ਲਈ ਇਕ ਪੇਟੀ ਬਣਾ ਰਿਹਾ ਹੈ। ਫਿਰ, ਕਿਉਂਕਿ ਕਾਫ਼ੀ ਦੇਰ ਹੋ ਚੁੱਕੀ ਸੀ, ਅਸੀਂ ਘਰ ਮੁੜ ਆਏ।

%%K t14-tit-3 sub
\VUpk{10.2pt}{40}{ਸ਼ਹਿਰ ਵਿਚ ਇਕ ਸੈਰ।}
\VUsaut{3.0mm}

%%K t14-07-1 p
7. --- ਅਗਲੇ ਦਿਨ ਮੇਰੀ ਮਾਂ ਨੇ ਸਲਾਹ ਦਿੱਤੀ ਕਿ ਜਾਨ ਆਪਣੀ ਤਸਵੀਰ ਖਿਚਵਾਏ, ਤੇ ਮੈਨੂੰ ਆਖਿਆ
ਕਿ ਉਸ ਨਾਲ \VUgras{ਤਸਵੀਰਸਾਜ਼} \textsuperscript{(20)} ਦੇ ਕੋਲ ਜਾਵਾਂ। ਖਾਣੇ ਪਿੱਛੋਂ ਅਸੀਂ
ਕਾਰੀਗਰ ਦੇ \VUgras{ਕਮਰੇ} ਵਿਚ ਗਏ, ਜਿੱਥੇ ਸਾਨੂੰ ਕਾਫ਼ੀ ਦੇਰ ਉਡੀਕ ਕਰਨੀ ਪਈ। ਵੇਲਾ ਕੱਟਣ ਨੂੰ
ਅਸੀਂ ਕੰਧ ਉੱਤੇ ਟੰਗੀਆਂ \VUgras{ਤਸਵੀਰਾਂ} \textsuperscript{(22)} ਵੇਖੀਆਂ।

%%K t14-07-2 p
ਸਾਡੀ ਵਾਰੀ ਆਉਣ ਉੱਤੇ ਕੰਮ ਵਿਚ ਦੇਰ ਨਾ ਲੱਗੀ, ਤੇ ਮੇਰੇ ਮਿੱਤਰ ਦੇ \VUgras{ਕੈਮਰੇ}
\textsuperscript{(21)} ਦੇ ਸਾਹਮਣੇ ਬਹਿ ਚੁੱਕਣ ਉੱਤੇ ਅਸੀਂ ਹੇਠਾਂ ਉਤਰ ਆਏ।

%%K t14-08-1 p
8. --- ਪਹਿਲੀ ਮੰਜ਼ਿਲ ਉੱਤੇ ਅਸੀਂ \VUgras{ਹਜਾਮ} \textsuperscript{(23)} ਦੇ ਖੁੱਲ੍ਹੇ ਬੂਹੇ
ਵਿਚੋਂ ਵੇਖਿਆ ਕਿ ਉਸ ਦਾ ਸਹਾਇਕ ਇਕ ਗਾਹਕ ਦੇ ਵਾਲ ਕੱਟ ਰਿਹਾ ਹੈ।

%%K t14-08-2 p
ਫਿਰ ਸੜਕ ਉੱਤੇ ਆ ਕੇ ਅਸੀਂ \VUgras{ਤੰਬਾਕੂ ਦੀ ਦੁਕਾਨ} \textsuperscript{(24)} ਕੋਲੋਂ
ਲੰਘੇ। ਜਾਨ ਨੂੰ ਲਾਲ \VUgras{ਲਾਲਟੈਣ} \textsuperscript{(25)} ਤੇ \VUgras{ਬੋਰਡ}
\textsuperscript{(26)} ਦੇ ਤੌਰ ਉੱਤੇ ਲੱਗੀ \VUgras{ਵੱਡੀ ਚਿਲਮ} ਤੋਂ ਇੰਨਾ ਮਜ਼ਾ ਆਇਆ ਕਿ ਉਹ
ਦੋ ਬੁੱਢੇ ਸੱਜਣਾਂ ਨਾਲ ਟਕਰਾ ਗਿਆ, ਜੋ \VUgras{ਨਸਵਾਰ ਦੀ ਚੁਟਕੀ} \textsuperscript{(27)}
ਲੈਂਦੇ ਗੱਲਾਂ ਕਰ ਰਹੇ ਸਨ।

%%K t14-tit-4 sub
\VUpk{10.2pt}{40}{ਮਿਠਾਈਆਂ ਦੀ ਦੁਕਾਨ।}
\VUsaut{3.0mm}

%%K t14-09-1 p
9. --- \VUgras{ਸਰਾਫ਼} \textsuperscript{(29)} ਦੀ \VUgras{ਖਿੜਕੀ} ਵਿਚ ਸਜੇ ਹਰ ਦੇਸ ਦੇ
\VUgras{ਨੋਟ} ਤੇ \VUgras{ਸਿੱਕੇ} ਨੇ ਕੁਝ ਦੇਰ ਸਾਡਾ ਧਿਆਨ ਬੰਨ੍ਹੀ ਰੱਖਿਆ, ਪਰ ਕਿਉਂਕਿ ਚਾਹ
ਦਾ ਵੇਲਾ ਸੀ, ਅਸੀਂ ਬਿਨਾਂ ਹੋਰ ਰੁਕੇ \VUgras{ਹਲਵਾਈ} \textsuperscript{(31)} ਦੀ ਦੁਕਾਨ ਵਿਚ
ਵੜ ਗਏ।

%%K t14-10-1 p
10. --- ਦੁਕਾਨ ਵਿਚ ਜਿੰਨੀਆਂ \VUgras{ਚੰਗੀਆਂ ਚੀਜ਼ਾਂ} ਸਨ — ਹਰ ਤਰ੍ਹਾਂ ਦੇ
\VUgras{ਕੇਕ, ਸੁੰਢ ਦੀ ਰੋਟੀ, ਬਿਸਕੁਟ} ਤੇ \VUgras{ਮਿਠਾਈਆਂ} — ਉਨ੍ਹਾਂ ਨੂੰ ਵੇਖ ਕੇ
ਸਾਥੋਂ ਚੋਣ ਹੀ ਨਾ ਹੋਈ। ਅਖ਼ੀਰ ਜਾਨ ਨੇ \VUgras{ਮਲਾਈ ਦੇ ਕੇਕ} ਲਏ, ਤੇ ਮੈਂ ਸਿਰਫ਼ ਇਕ ਨਿੱਕੀ
\VUgras{ਚੈਰੀ ਦੀ ਟਿੱਕੀ}, ਕਿਉਂਕਿ ਮਿੱਠੀਆਂ ਚੀਜ਼ਾਂ ਨਾਲ \VUgras{ਮੇਰੇ ਦੰਦ ਦੁਖਦੇ ਹਨ} ਤੇ
ਮੈਨੂੰ \VUgras{ਦੰਦਾਂ ਦੇ ਡਾਕਟਰ} \textsuperscript{(30)} ਦੇ ਕੋਲ ਵਾਰ ਵਾਰ ਜਾਣ ਦੀ ਰਤਾ
ਇੱਛਾ ਨਹੀਂ।

%%K t14-tit-5 sub
\VUpk{10.2pt}{40}{ਟੰਕਣ।}
\VUsaut{3.0mm}

%%K t14-11-1 p
11. --- ਆਪਣੇ ਮਿੱਤਰ ਨੂੰ ਇਕ \VUgras{ਟਾਈਪਰਾਈਟਰ} ਵਿਖਾਉਣ ਲਈ, ਜਿਸ ਬਾਰੇ ਉਸ ਨੇ ਸੁਣਿਆ ਤਾਂ
ਸੀ ਪਰ ਵੇਖਿਆ ਨਾ ਸੀ, ਅਸੀਂ ਉਸ \VUgras{ਦਫ਼ਤਰ} \textsuperscript{(32)} ਵਿਚ ਗਏ ਜੋ ਮੇਰੇ
\VUgras{ਚਚੇਰੇ ਭਰਾ} \textsuperscript{(34)} ਦਾ ਹੈ। ਅਸੀਂ ਉਨ੍ਹਾਂ ਨੂੰ ਆਪਣੀ
\VUgras{ਸਕੱਤਰ} \textsuperscript{(35)} ਨੂੰ ਖ਼ਤ ਲਿਖਵਾਉਂਦੇ ਲੱਭਿਆ, ਜੋ
\VUgras{ਆਸ਼ੂਲਿਪਿਕ} ਹੈ। ਇਕ ਖ਼ਤ ਮੁੱਕਦੇ ਹੀ ਇਕ \VUgras{ਟੰਕਕ} \textsuperscript{(33)} ਉਸ
ਨੂੰ ਆਪਣੀ ਮਸ਼ੀਨ ਉੱਤੇ ਉਤਾਰ ਲੈਂਦਾ ਸੀ। ਆਪਣੇ ਰਿਸ਼ਤੇਦਾਰ ਦਾ ਕੰਮ ਨਾ ਰੋਕਣਾ ਚਾਹੁੰਦੇ ਹੋਏ ਅਸੀਂ
ਉਨ੍ਹਾਂ ਕੋਲ ਸਿਰਫ਼ ਕੁਝ ਮਿੰਟ ਰੁਕੇ।

%%K t14-12-1 p
12. --- ਮੇਰੇ ਚਚੇਰੇ ਭਰਾ ਦੇ ਦਫ਼ਤਰ ਦੇ ਹੇਠਾਂ ਇਕ ਵੱਡੇ \VUgras{ਘੜੀਸਾਜ਼ ਤੇ ਜੌਹਰੀ}
\textsuperscript{(37)} ਰਹਿੰਦੇ ਹਨ, ਜੋ ਚਾਂਦੀ ਦੇ ਕਾਰੀਗਰ ਵੀ ਹਨ। ਉਨ੍ਹਾਂ ਦੀ ਸ਼ਾਨਦਾਰ ਦੁਕਾਨ
ਵਿਚ ਬਹੁਤ ਸਾਰੀਆਂ \VUgras{ਘੜੀਆਂ, ਅੰਗੀਠੀ ਦੀਆਂ ਘੜੀਆਂ, ਜੇਬ-ਘੜੀਆਂ},
\VUgras{ਜ਼ੰਜੀਰਾਂ} ਤੇ ਕੀਮਤੀ \VUgras{ਗਹਿਣੇ} ਹਨ, ਜਿਵੇਂ \VUgras{ਮੋਤੀਆਂ ਦੇ ਹਾਰ},
ਸੋਨੇ ਦੇ \VUgras{ਕੰਗਣ}, \VUgras{ਕੰਨ-ਫੁੱਲ} ਤੇ \VUgras{ਕੀਮਤੀ ਪੱਥਰਾਂ} ਤੇ
\VUgras{ਹੀਰਿਆਂ} ਨਾਲ ਜੜੀਆਂ \VUgras{ਮੁੰਦਰੀਆਂ}। ਇਕ ਖਿੜਕੀ ਪੂਰੀ ਤਰ੍ਹਾਂ
\VUgras{ਚਾਂਦੀ ਦੇ ਭਾਂਡਿਆਂ} ਨੂੰ ਦਿੱਤੀ ਗਈ ਹੈ ਤੇ ਉਸ ਵਿਚ ਚਾਂਦੀ ਦੇ
\VUgras{ਚਮਚਿਆਂ-ਕਾਂਟਿਆਂ} ਦਾ ਇਕ ਵੱਡਾ ਸੰਗ੍ਰਹਿ ਹੈ।

%%K t14-13-1 p
13. --- ਸੜਕ ਦਾ ਮੋੜ ਮੁੜਦੇ ਹੋਏ ਮੈਂ ਵੇਖਿਆ ਕਿ ਘਰ ਦੇ \VUgras{ਨੰਬਰ} ਕੋਲ ਲੱਗੀ
\VUgras{ਨਾਂ ਦੀ ਪੱਟੀ} \textsuperscript{(36)} ਬਦਲ ਦਿੱਤੀ ਗਈ ਹੈ। ਮੈਂ ਇਹ ਜ਼ੋਰ ਨਾਲ ਆਖਣ ਹੀ
ਵਾਲਾ ਸੀ ਕਿ ਇਕ ਫ਼ੌਜੀ ਵਾਜੇ ਦੀ ਖ਼ੁਸ਼ ਆਵਾਜ਼ ਸੁਣਾਈ ਦਿੱਤੀ। ਅਸੀਂ ਰਜਮੰਟ ਵੱਲ ਦੌੜ ਪਏ, ਇਹ ਸੋਚੇ
ਬਿਨਾਂ ਕਿ ਉਹ \VUgras{ਟੋਪੀ ਵਾਲੇ} \textsuperscript{(39)} ਤੇ \VUgras{ਦਸਤਾਨੇ ਵਾਲੇ}
\textsuperscript{(40)} ਦੀਆਂ ਦੁਕਾਨਾਂ ਦੇ ਸਾਹਮਣੇ ਲੰਘਣ ਵਾਲੀ ਹੈ, ਤੇ ਕਿ ਉਸ ਨੂੰ ਲੰਘਦੇ ਵੇਖਣ
ਲਈ ਅਸੀਂ \VUgras{ਐਨਕ ਵਾਲੇ} \textsuperscript{(38)} ਦੀ ਖਿੜਕੀ ਦੇ ਸਾਹਮਣੇ ਖੜੇ ਰਹਿ ਕੇ ਵੀ
ਓਨੀ ਹੀ ਚੰਗੀ ਥਾਂ ਉੱਤੇ ਹੁੰਦੇ, ਜੋ \VUgras{ਨੱਕ-ਐਨਕਾਂ, ਐਨਕਾਂ} ਤੇ \VUgras{ਦੂਰਬੀਨਾਂ}
ਨਾਲ ਭਰੀ ਹੈ।

%%K t14-tit-6 sub
\VUpk{10.2pt}{40}{ਸਲਾਮੀ ਪਰੇਡ।}
\VUsaut{3.0mm}

%%K t14-14-1 p
14. --- \VUgras{ਰਜਮੰਟ} \textsuperscript{(41)} ਦੇ ਅੱਗੇ ਅੱਗੇ \VUgras{ਢੋਲ-ਮੁਖੀ}
\textsuperscript{(42)} ਤੁਰ ਰਿਹਾ ਸੀ, ਉਸ ਦੇ ਪਿੱਛੇ \VUgras{ਢੋਲ ਵਾਲੇ}
\textsuperscript{(43)}, \VUgras{ਬਿਗਲ ਵਾਲੇ} \textsuperscript{(44)} ਤੇ ਸਾਰਾ
\VUgras{ਵਾਜਾ}। \VUgras{ਜਰਨੈਲ} \textsuperscript{(45)}, ਜੋ ਆਪਣੇ ਸਹਾਇਕ ਨਾਲ ਚੌਕ ਉੱਤੇ
ਸਨ, \VUgras{ਪਾਣੀ ਦੀ ਧਾਰ} \textsuperscript{(49)} ਕੋਲ ਰੁਕ ਗਏ ਸਨ, ਜੋ
\VUgras{ਫੁਹਾਰੇ} \textsuperscript{(48)} ਦੀ ਹੈ, \VUgras{ਪਰੇਡ} ਵੇਖਣ ਲਈ।

%%K t14-14-2 p
\VUgras{ਅਮਲੇ ਦੇ ਅਫ਼ਸਰਾਂ} \textsuperscript{(46)}, \VUgras{ਕਰਨਲ} ਤੇ
\VUgras{ਨਾਇਬ ਕਰਨਲ} ਨੇ ਉਨ੍ਹਾਂ ਨੂੰ ਤਲਵਾਰ ਨਾਲ ਸਲਾਮੀ ਦਿੱਤੀ। \VUgras{ਮੇਜਰ} ਆਪਣੀਆਂ
\VUgras{ਬਟਾਲੀਅਨਾਂ} \textsuperscript{(47)} ਦੇ ਅੱਗੇ ਸਨ ਤੇ ਹਰ \VUgras{ਕਪਤਾਨ} ਆਪਣੀ
\VUgras{ਕੰਪਨੀ} ਦੀ ਅਗਵਾਈ ਕਰਦਾ ਸੀ, ਜਦ ਕਿ \VUgras{ਲੈਫ਼ਟੀਨੈਂਟ},
\VUgras{ਨਾਇਬ ਲੈਫ਼ਟੀਨੈਂਟ} ਤੇ \VUgras{ਗ਼ੈਰ-ਕਮਿਸ਼ਨ ਅਫ਼ਸਰ} ਆਪਣੇ ਬੰਦਿਆਂ ਕੋਲ ਆਪਣੀਆਂ ਆਮ
ਥਾਵਾਂ ਉੱਤੇ ਸਨ।

%%K t14-15-1 p
15. --- \VUgras{ਚੁਰਾਹੇ} \textsuperscript{(50)} ਉੱਤੇ ਬਹੁਤ ਸਾਰੇ ਰਾਹੀ ਜਮ੍ਹਾਂ ਹੋ ਗਏ ਸਨ;
ਦੁਕਾਨਦਾਰ — \VUgras{ਛਤਰੀ ਵੇਚਣ ਵਾਲਾ} \textsuperscript{(51)}, \VUgras{ਰੰਗਰੇਜ਼}
\textsuperscript{(53)}, \VUgras{ਬੰਦੂਕਸਾਜ਼} \textsuperscript{(54)} — \VUgras{ਸਿਪਾਹੀ}
ਵੇਖਣ ਬਾਹਰ ਆ ਗਏ। ਹਰ ਗੱਡੀ — \VUgras{ਕਿਰਾਏ ਦੀਆਂ ਬੱਘੀਆਂ} \textsuperscript{(52)},
\VUgras{ਨਿੱਜੀ ਗੱਡੀਆਂ} \textsuperscript{(57)} ਤੇ \VUgras{ਪਾਣੀ ਦੀ ਗੱਡੀ}
\textsuperscript{(56)} ਵੀ — ਪਟੜੀ ਨਾਲ ਖੜੀ ਹੋ ਗਈ।

%%K t14-tit-7 sub
\VUpk{10.2pt}{40}{ਸੜਕ ਦੇ ਦ੍ਰਿਸ਼।}
\VUsaut{3.0mm}

%%K t14-15-2 p
ਹੱਥ ਵਿਚ ਆਪਣੇ \VUgras{ਨਮੂਨੇ ਦੇ ਬਕਸੇ} ਚੁੱਕੀ ਇਕ \VUgras{ਸਫ਼ਰੀ ਵੇਚਣ ਵਾਲਾ}
\textsuperscript{(55)} ਫੁਰਤੀ ਨਾਲ ਸੜਕ ਲੰਘ ਗਿਆ, ਤੇ \VUgras{ਉੱਪਰਲੇ ਪਟੜੇ}
\textsuperscript{(59)} ਉੱਤੇ ਬੈਠੇ ਲੋਕ, ਜੋ \VUgras{ਓਮਨੀਬੱਸ} \textsuperscript{(58)} ਦਾ
ਹੈ, ਮੁੜ ਕੇ ਤਮਾਸ਼ਾ ਵੇਖਣ ਲੱਗੇ। ਇਕ ਵਿਚਾਰੇ \VUgras{ਗਲੀ ਦੇ ਗਵੱਈਏ}
\textsuperscript{(60)} ਨੂੰ, ਜੋ ਆਪਣਾ \VUgras{ਪੇਟੀ-ਵਾਜਾ} \textsuperscript{(61)} ਚੁੱਕੀ
ਸੀ, ਆਪਣਾ \VUgras{ਗੀਤ} ਵਿਚਕਾਰ ਛੱਡਣਾ ਪਿਆ, ਕਿਉਂਕਿ ਢੋਲਾਂ ਦੇ ਰੌਲੇ ਨੇ ਉਸ ਦੀ ਆਵਾਜ਼ ਦੱਬ
ਦਿੱਤੀ।

%%K t14-16-1 p
16. --- ਜਦੋਂ ਰਜਮੰਟ \VUgras{ਕੱਪੜੇ ਦੀ ਦੁਕਾਨ} \textsuperscript{(64)} ਦੇ ਸਾਹਮਣੇ ਆਈ, ਤਾਂ
\VUgras{ਸਿਲਾਈ ਵਾਲੀਆਂ} \textsuperscript{(63)} ਤੇ \VUgras{ਦਰਜ਼ੀ}
\textsuperscript{(62)} ਆਪਣੀਆਂ \VUgras{ਸਿਲਾਈ ਦੀਆਂ ਮਸ਼ੀਨਾਂ} ਤੇ ਆਪਣਾ ਕੰਮ ਛੱਡ ਕੇ ਖਿੜਕੀ
ਵਿਚੋਂ ਝਾਕਣ ਲੱਗੇ। \VUgras{ਦੁਕਾਨਦਾਰ} \textsuperscript{(65)}, ਜਿਸ ਨੇ ਹੁਣੇ ਹੁਣੇ ਆਪਣੀ
\VUgras{ਸਜਾਵਟ ਦੀ ਖਿੜਕੀ} ਵਿਚ ਨਵੇਂ \VUgras{ਸੂਟ} \textsuperscript{(66)} ਰੱਖੇ ਸਨ, ਉਸ
ਨੂੰ \VUgras{ਕੱਪੜੇ} \textsuperscript{(67)} ਤੇ \VUgras{ਲਿਨਨ} ਦੇ ਉਹ ਥਾਨ ਹਟਾਉਣੇ ਪਏ ਜੋ
ਪਟੜੀ ਉੱਤੇ \VUgras{ਨਾਲੀ ਦੇ ਮੂੰਹ} \textsuperscript{(68)} ਕੋਲ ਰੱਖੇ ਸਨ, ਇਸ ਡਰ ਤੋਂ ਕਿ ਭੀੜ
ਉਨ੍ਹਾਂ ਨੂੰ ਡੇਗ ਨਾ ਦੇਵੇ; ਇਕ ਬੁੱਢੇ \VUgras{ਫੇਰੀ ਵਾਲੇ} \textsuperscript{(69)} ਤਕ ਨੂੰ,
ਜਿਸ ਨਾਲ ਇਕ ਚੌਕੀਦਾਰਨੀ ਕੁਝ ਜੁਰਾਬਾਂ ਦਾ ਮੁੱਲ-ਭਾਅ ਕਰ ਰਹੀ ਸੀ, ਆਪਣੀ ਵਿਕਰੀ ਪੂਰੀ ਕਰਨ ਲਈ ਇਕ
ਡਿਉੜ੍ਹੀ ਵਿਚ ਵੜਨਾ ਪਿਆ।

%%K t14-16-2 p
ਅਸੀਂ ਰਜਮੰਟ ਨਾਲ ਬੈਰਕਾਂ ਤਕ ਗਏ \VUgras{ਤੇ}, ਆਪਣੇ ਦਿਨ ਤੋਂ ਬਹੁਤ ਤਸੱਲੀ ਪਾ ਕੇ, ਖਾਣੇ ਦੇ
ਵੇਲੇ ਘਰ ਮੁੜ ਆਏ।

%%K t14-tit-8 sub
\VUpk{10.2pt}{40}{ਵੱਖੋ-ਵੱਖ ਵਪਾਰ ਤੇ ਪੇਸ਼ੇ।}
\VUsaut{3.0mm}

%%K t14-17-1 p
17. --- ਅਗਲੇ ਦਿਨ ਮੇਰੇ ਪਿਤਾ ਨੇ ਸਲਾਹ ਦਿੱਤੀ ਕਿ ਸਾਨੂੰ ਸਥਾਨਕ ਅਖ਼ਬਾਰ ਦਾ ਛਾਪਾਖ਼ਾਨਾ ਵਿਖਾ
ਲਿਆਉਣ। ਅਸੀਂ ਜੋਸ਼ ਨਾਲ ਮੰਨ ਲਿਆ।

%%K t14-17-2 p
ਉਹ \VUgras{ਛਾਪੇ ਵਾਲਾ} \VUgras{ਕਿਤਾਬਾਂ ਦਾ ਵੇਚਣ ਵਾਲਾ} \textsuperscript{(70)},
\VUgras{ਪ੍ਰਕਾਸ਼ਕ}, \VUgras{ਲਿਖਣ ਦੇ ਸਾਮਾਨ ਦਾ ਵੇਚਣ ਵਾਲਾ} ਤੇ \VUgras{ਜਿਲਦਸਾਜ਼} ਵੀ
ਹੈ। ਉਸ ਨੇ ਪਹਿਲੀ ਮੰਜ਼ਿਲ ਉੱਤੇ ਅਖ਼ਬਾਰ ਦਾ \VUgras{ਸੰਪਾਦਕੀ ਕਮਰਾ} \textsuperscript{(71)}
ਰੱਖਿਆ ਹੈ, ਤੇ \VUgras{ਨੱਕਾਸ਼ੀ ਦਾ ਕਮਰਾ} ਵੀ, ਜਿੱਥੇ \VUgras{ਪੱਥਰ-ਛਾਪਕ}
\textsuperscript{(72)} ਤੇ \VUgras{ਤਾਂਬੇ ਦੇ ਨੱਕਾਸ਼} \textsuperscript{(73)} ਕੰਮ ਕਰਦੇ
ਹਨ, ਤੇ ਅੰਤ ਵਿਚ ਜੋੜਨ ਦਾ ਕਮਰਾ, ਜਿੱਥੇ \VUgras{ਅੱਖਰ ਜੋੜਨ ਵਾਲੇ} \textsuperscript{(74)}
ਹਨ। ਅਸਲੀ \VUgras{ਛਾਪਾਖ਼ਾਨਾ} \textsuperscript{(75)}, \VUgras{ਛਾਪਣ ਦੀਆਂ ਮਸ਼ੀਨਾਂ} ਤੇ
ਵੱਖੋ-ਵੱਖ ਜੰਤਰ ਹੇਠਲੀ ਮੰਜ਼ਿਲ ਉੱਤੇ ਹਨ।

%%K t14-tit-9 sub
\VUpk{10.2pt}{40}{ਜਾਨ ਬਹੁਤ ਸਾਰੇ ਸਵਾਲ ਪੁੱਛਦਾ ਹੈ।}
\VUsaut{3.0mm}

%%K t14-18-1 p
18. --- ਛਾਪਾਖ਼ਾਨੇ ਤੋਂ ਨਿਕਲਦੇ ਹੋਏ ਜਾਨ ਬੜੀ ਉਲਝਣ ਵਿਚ ਇਕ ਥੰਮ੍ਹ ਉੱਤੇ ਲੱਗੇ ਨਿੱਕੇ ਸ਼ੀਸ਼ੇ
ਦੇ ਡੱਬੇ ਦੇ ਸਾਹਮਣੇ ਰੁਕ ਗਿਆ, ਜੋ \VUgras{ਕਿਰਿਆਨੇ ਵਾਲੇ} \textsuperscript{(77)} ਦੇ
ਸਾਹਮਣੇ ਸੀ, ਤੇ ਪੁੱਛਿਆ ਕਿ ਉਹ ਕਿਸ ਕੰਮ ਦਾ ਹੈ। ਮੇਰੇ ਪਿਤਾ ਨੇ ਸਮਝਾਇਆ ਕਿ ਉਹ
\VUgras{ਅੱਗ ਦੀ ਸੂਚਨਾ} \textsuperscript{(76)} ਹੈ। ਸੱਚਮੁੱਚ ਸ਼ਹਿਰ ਦੀ ਹਰ ਚੀਜ਼ ਮੇਰੇ ਮਿੱਤਰ
ਲਈ ਅਚਰਜ ਸੀ, ਸਧਾਰਨ \VUgras{ਪਾਣੀ ਦੇ ਨਲ} \textsuperscript{(78)} ਤੇ ਉਸ ਹੌਲੇ
\VUgras{ਸਾਮਾਨ ਦੇ ਤਿਪਹੀਏ} \textsuperscript{(80)} ਤੋਂ ਲੈ ਕੇ, ਜਿਸ ਨੂੰ ਵਰਦੀ ਪਾਈ ਇਕ
\VUgras{ਛੋਕਰਾ} \textsuperscript{(79)} ਚਲਾਉਂਦਾ ਹੈ, \VUgras{ਡਾਕ ਦੀ ਗੱਡੀ}
\textsuperscript{(81)} ਤਕ।

%%K t14-19-1 p
19. --- ਜਦੋਂ ਤਕ ਮੇਰੇ ਪਿਤਾ ਸਾਨੂੰ ਛੱਡ ਕੇ \VUgras{ਮਹਿਸੂਲ ਵਸੂਲਣ ਵਾਲੇ}
\textsuperscript{(83)} ਨਾਲ ਗੱਲ ਕਰਨ ਗਏ, ਜੋ ਇਕ \VUgras{ਵਕੀਲ} \textsuperscript{(82)}
ਤੇ ਇਕ \VUgras{ਮੁਨਸ਼ੀ} \textsuperscript{(84)} ਨਾਲ ਗੱਲਾਂ ਕਰਨ ਰੁਕੇ ਸਨ, ਅਸੀਂ ਸੜਕ ਦੀ
ਚਹਿਲ-ਪਹਿਲ ਵੇਖਣ ਵਿਚ ਦਿਲ ਪਰਚਾਉਂਦੇ ਰਹੇ। ਸਾਡੇ ਸਾਹਮਣੇ ਤੋਂ ਕਈ ਤਰ੍ਹਾਂ ਦੇ ਲੋਕ ਲੰਘੇ : ਇਕ
\VUgras{ਹਲਵਾਈ ਦਾ ਛੋਕਰਾ} \textsuperscript{(85)}, ਜੋ ਇਕ ਟੋਕਰੀ ਵਿਚ ਇਕ ਸ਼ਾਨਦਾਰ
\VUgras{ਪਾਈ} ਚੁੱਕੀ ਸੀ, ਇਕ \VUgras{ਕੱਪੜਿਆਂ ਦੇ ਵਪਾਰੀ} \textsuperscript{(86)} ਦੇ
ਪਿੱਛੇ ਆਇਆ; ਇਕ \VUgras{ਬੱਤੀ ਬਾਲਣ ਵਾਲਾ} \textsuperscript{(87)} ਇਕ
\VUgras{ਗੈਸ ਲਾਉਣ ਵਾਲੇ} \textsuperscript{(88)} ਨਾਲ ਗੱਲਾਂ ਕਰ ਰਿਹਾ ਸੀ; ਇਕ
\VUgras{ਸਾਧਵੀ} \textsuperscript{(89)} ਇਕ ਨਿੱਕੀ ਯਤੀਮ ਕੁੜੀ ਦਾ ਹੱਥ ਫੜੀ ਆਪਣੇ ਮੱਠ ਨੂੰ ਮੁੜ
ਰਹੀ ਸੀ।

%%K t14-tit-10 sub
\VUpk{10.2pt}{40}{ਮੁੜਦੇ ਵੇਲੇ।}
\VUsaut{3.0mm}

%%K t14-20-1 p
20. --- ਮੇਰੇ ਪਿਤਾ ਸਾਡੇ ਨਾਲ ਆ ਰਲੇ ਹੀ ਸਨ ਕਿ ਜਾਨ ਦੋ \VUgras{ਚਿਮਨੀ ਸਾਫ਼ ਕਰਨ ਵਾਲਿਆਂ}
\textsuperscript{(91)} ਦੇ ਕਾਲਖ਼ ਨਾਲ ਕਾਲੇ ਮੂੰਹਾਂ ਤੇ ਅਨੋਖੇ ਭੇਸ ਉੱਤੇ ਹੱਸ ਪਿਆ।

%%K t14-21-1 p
21. --- ਮੈਂ ਆਪਣੀ ਮਾਂ ਲਈ \VUgras{ਫੁੱਲ ਵੇਚਣ ਵਾਲੀ} \textsuperscript{(92)} ਕੋਲੋਂ ਇਕ
ਬੂਟਾ ਖ਼ਰੀਦਣਾ ਚਾਹੁੰਦਾ ਸੀ, ਪਰ ਮੇਰੇ ਪਿਤਾ ਨੇ ਆਖਿਆ ਕਿ ਉਹ ਢੋਣ ਵਿਚ ਬਹੁਤ ਭਾਰਾ ਹੋਵੇਗਾ ਤੇ
ਕੁਝ \VUgras{ਸੰਤਰੇ} ਲੈ ਲੈਣੇ ਚੰਗੇ ਰਹਿਣਗੇ। ਅਸੀਂ \VUgras{ਰੇੜ੍ਹੀ ਵਾਲੀ}
\textsuperscript{(94)} ਕੋਲ ਗਏ, ਤੇ ਜਦੋਂ \VUgras{ਅਧਿਆਪਕਾ} \textsuperscript{(93)}, ਜੋ
ਆਪਣੀ ਨਿੱਕੀ ਸ਼ਾਗਿਰਦ ਲਈ ਕੁਝ ਛਾਂਟ ਰਹੀ ਸੀ, ਆਪਣੀ ਖ਼ਰੀਦ ਪੂਰੀ ਕਰ ਚੁੱਕੀ, ਤਦੋਂ ਸਾਡੀ ਵਾਰੀ
ਆਈ।

%%K t14-22-1 p
22. --- ਮੁੜਦੇ ਹੋਏ ਸਾਨੂੰ ਕਈ ਜਾਣ-ਪਛਾਣ ਵਾਲੇ ਮਿਲੇ : ਮੇਰੇ ਪਰਿਵਾਰ ਦੀ ਇਕ ਪੁਰਾਣੀ ਸਹੇਲੀ,
ਆਪਣੀ \VUgras{ਸੰਗਣੀ} \textsuperscript{(95)} ਦੀ ਬਾਂਹ ਫੜੀ, ਪੁਲਾਂ ਤੇ ਸੜਕਾਂ ਦੇ
\VUgras{ਇੰਜੀਨੀਅਰ} \textsuperscript{(96)}, ਤੇ ਮੇਰੀ ਇਕ ਚਚੇਰੀ ਭੈਣ ਆਪਣੀ \VUgras{ਆਇਆ}
\textsuperscript{(98)} ਨਾਲ।

%%K t14-22-2 p
ਫਿਰ ਅਸੀਂ ਮੇਰੀ ਮਾਂ ਦੀ \VUgras{ਟੋਪੀਆਂ ਬਣਾਉਣ ਵਾਲੀ} \textsuperscript{(97)} ਕੋਲੋਂ ਲੰਘੇ
ਤੇ ਦੋ \VUgras{ਸਾਧੂਆਂ} \textsuperscript{(99)} ਕੋਲੋਂ, ਜੋ ਸ਼ਹਿਰ ਵਿਚ ਓਪਰੇ ਸਨ ਤੇ ਜੋ ਮੇਰੇ
ਪਿਤਾ ਕੋਲੋਂ ਸਟੇਸ਼ਨ ਦਾ ਰਾਹ ਪੁੱਛਣ ਆਏ।
\VUsaut{6.0mm}
\VUfilet{22mm}
